% Options for packages loaded elsewhere
\PassOptionsToPackage{unicode}{hyperref}
\PassOptionsToPackage{hyphens}{url}
\documentclass[
]{article}
\usepackage{xcolor}
\usepackage[margin=1in]{geometry}
\usepackage{amsmath,amssymb}
\setcounter{secnumdepth}{-\maxdimen} % remove section numbering
\usepackage{iftex}
\ifPDFTeX
  \usepackage[T1]{fontenc}
  \usepackage[utf8]{inputenc}
  \usepackage{textcomp} % provide euro and other symbols
\else % if luatex or xetex
  \usepackage{unicode-math} % this also loads fontspec
  \defaultfontfeatures{Scale=MatchLowercase}
  \defaultfontfeatures[\rmfamily]{Ligatures=TeX,Scale=1}
\fi
\usepackage{lmodern}
\ifPDFTeX\else
  % xetex/luatex font selection
\fi
% Use upquote if available, for straight quotes in verbatim environments
\IfFileExists{upquote.sty}{\usepackage{upquote}}{}
\IfFileExists{microtype.sty}{% use microtype if available
  \usepackage[]{microtype}
  \UseMicrotypeSet[protrusion]{basicmath} % disable protrusion for tt fonts
}{}
\makeatletter
\@ifundefined{KOMAClassName}{% if non-KOMA class
  \IfFileExists{parskip.sty}{%
    \usepackage{parskip}
  }{% else
    \setlength{\parindent}{0pt}
    \setlength{\parskip}{6pt plus 2pt minus 1pt}}
}{% if KOMA class
  \KOMAoptions{parskip=half}}
\makeatother
\usepackage{color}
\usepackage{fancyvrb}
\newcommand{\VerbBar}{|}
\newcommand{\VERB}{\Verb[commandchars=\\\{\}]}
\DefineVerbatimEnvironment{Highlighting}{Verbatim}{commandchars=\\\{\}}
% Add ',fontsize=\small' for more characters per line
\usepackage{framed}
\definecolor{shadecolor}{RGB}{248,248,248}
\newenvironment{Shaded}{\begin{snugshade}}{\end{snugshade}}
\newcommand{\AlertTok}[1]{\textcolor[rgb]{0.94,0.16,0.16}{#1}}
\newcommand{\AnnotationTok}[1]{\textcolor[rgb]{0.56,0.35,0.01}{\textbf{\textit{#1}}}}
\newcommand{\AttributeTok}[1]{\textcolor[rgb]{0.13,0.29,0.53}{#1}}
\newcommand{\BaseNTok}[1]{\textcolor[rgb]{0.00,0.00,0.81}{#1}}
\newcommand{\BuiltInTok}[1]{#1}
\newcommand{\CharTok}[1]{\textcolor[rgb]{0.31,0.60,0.02}{#1}}
\newcommand{\CommentTok}[1]{\textcolor[rgb]{0.56,0.35,0.01}{\textit{#1}}}
\newcommand{\CommentVarTok}[1]{\textcolor[rgb]{0.56,0.35,0.01}{\textbf{\textit{#1}}}}
\newcommand{\ConstantTok}[1]{\textcolor[rgb]{0.56,0.35,0.01}{#1}}
\newcommand{\ControlFlowTok}[1]{\textcolor[rgb]{0.13,0.29,0.53}{\textbf{#1}}}
\newcommand{\DataTypeTok}[1]{\textcolor[rgb]{0.13,0.29,0.53}{#1}}
\newcommand{\DecValTok}[1]{\textcolor[rgb]{0.00,0.00,0.81}{#1}}
\newcommand{\DocumentationTok}[1]{\textcolor[rgb]{0.56,0.35,0.01}{\textbf{\textit{#1}}}}
\newcommand{\ErrorTok}[1]{\textcolor[rgb]{0.64,0.00,0.00}{\textbf{#1}}}
\newcommand{\ExtensionTok}[1]{#1}
\newcommand{\FloatTok}[1]{\textcolor[rgb]{0.00,0.00,0.81}{#1}}
\newcommand{\FunctionTok}[1]{\textcolor[rgb]{0.13,0.29,0.53}{\textbf{#1}}}
\newcommand{\ImportTok}[1]{#1}
\newcommand{\InformationTok}[1]{\textcolor[rgb]{0.56,0.35,0.01}{\textbf{\textit{#1}}}}
\newcommand{\KeywordTok}[1]{\textcolor[rgb]{0.13,0.29,0.53}{\textbf{#1}}}
\newcommand{\NormalTok}[1]{#1}
\newcommand{\OperatorTok}[1]{\textcolor[rgb]{0.81,0.36,0.00}{\textbf{#1}}}
\newcommand{\OtherTok}[1]{\textcolor[rgb]{0.56,0.35,0.01}{#1}}
\newcommand{\PreprocessorTok}[1]{\textcolor[rgb]{0.56,0.35,0.01}{\textit{#1}}}
\newcommand{\RegionMarkerTok}[1]{#1}
\newcommand{\SpecialCharTok}[1]{\textcolor[rgb]{0.81,0.36,0.00}{\textbf{#1}}}
\newcommand{\SpecialStringTok}[1]{\textcolor[rgb]{0.31,0.60,0.02}{#1}}
\newcommand{\StringTok}[1]{\textcolor[rgb]{0.31,0.60,0.02}{#1}}
\newcommand{\VariableTok}[1]{\textcolor[rgb]{0.00,0.00,0.00}{#1}}
\newcommand{\VerbatimStringTok}[1]{\textcolor[rgb]{0.31,0.60,0.02}{#1}}
\newcommand{\WarningTok}[1]{\textcolor[rgb]{0.56,0.35,0.01}{\textbf{\textit{#1}}}}
\usepackage{graphicx}
\makeatletter
\newsavebox\pandoc@box
\newcommand*\pandocbounded[1]{% scales image to fit in text height/width
  \sbox\pandoc@box{#1}%
  \Gscale@div\@tempa{\textheight}{\dimexpr\ht\pandoc@box+\dp\pandoc@box\relax}%
  \Gscale@div\@tempb{\linewidth}{\wd\pandoc@box}%
  \ifdim\@tempb\p@<\@tempa\p@\let\@tempa\@tempb\fi% select the smaller of both
  \ifdim\@tempa\p@<\p@\scalebox{\@tempa}{\usebox\pandoc@box}%
  \else\usebox{\pandoc@box}%
  \fi%
}
% Set default figure placement to htbp
\def\fps@figure{htbp}
\makeatother
\setlength{\emergencystretch}{3em} % prevent overfull lines
\providecommand{\tightlist}{%
  \setlength{\itemsep}{0pt}\setlength{\parskip}{0pt}}
\usepackage{bookmark}
\IfFileExists{xurl.sty}{\usepackage{xurl}}{} % add URL line breaks if available
\urlstyle{same}
\hypersetup{
  pdftitle={Data\_Analysis\_Thesis},
  pdfauthor={Monique White},
  hidelinks,
  pdfcreator={LaTeX via pandoc}}

\title{Data\_Analysis\_Thesis}
\author{Monique White}
\date{2025-12-03}

\begin{document}
\maketitle

Clean data

\begin{Shaded}
\begin{Highlighting}[]
\NormalTok{Coral\_settlement\_nGBR\_CS }\OtherTok{\textless{}{-}} \FunctionTok{read\_csv}\NormalTok{(}\StringTok{"\textasciitilde{}/Larval recruitment resources/Data Analysis/Coral{-}Settlement{-}Northern{-}GBR/Coral settlement\_nGBR\_CS.csv"}\NormalTok{) }\SpecialCharTok{\%\textgreater{}\%}
  \FunctionTok{mutate}\NormalTok{(}\FunctionTok{across}\NormalTok{(}\DecValTok{15}\SpecialCharTok{:}\DecValTok{29}\NormalTok{, }\SpecialCharTok{\textasciitilde{}}\FunctionTok{replace\_na}\NormalTok{(.x, }\DecValTok{0}\NormalTok{))) }\CommentTok{\# Read in data}

\NormalTok{Coral\_settlement\_nGBR }\OtherTok{\textless{}{-}}\NormalTok{ Coral\_settlement\_nGBR\_CS }\SpecialCharTok{\%\textgreater{}\%} 
  \FunctionTok{filter}\NormalTok{(}\StringTok{\textasciigrave{}}\AttributeTok{Shelf posn}\StringTok{\textasciigrave{}} \SpecialCharTok{!=} \StringTok{"Coral Sea"}\NormalTok{, }\CommentTok{\#remove coral sea data}
        \SpecialCharTok{!}\NormalTok{Obs }\SpecialCharTok{\%in\%} \FunctionTok{c}\NormalTok{(}\DecValTok{1499}\NormalTok{, }\DecValTok{1621}\NormalTok{, }\DecValTok{1622}\NormalTok{)) }\CommentTok{\#Remove loose tiles with heavy biofouling}

\NormalTok{Coral\_settlement\_nGBR }\SpecialCharTok{|\textgreater{}} 
  \FunctionTok{filter}\NormalTok{(notes }\SpecialCharTok{!=} \StringTok{"NA"}\NormalTok{)}
\end{Highlighting}
\end{Shaded}

\begin{verbatim}
## # A tibble: 54 x 30
##      Obs `Year (Jan)` Region `Shelf posn` Reef     Site   `Site #` Habitat Tile 
##    <dbl>        <dbl> <chr>  <chr>        <chr>    <chr>  <chr>    <chr>   <chr>
##  1   625         2022 GBRMP  Inner        Turtle N TN_co~ 1        Crest   10   
##  2   627         2022 GBRMP  Inner        Turtle N TN_co~ 1        Crest   12   
##  3   631         2022 GBRMP  Inner        Turtle N TN_Ba~ 2        Crest   2    
##  4   658         2022 GBRMP  Inner        Turtle S TS_Cr~ 6        Crest   14   
##  5   661         2022 GBRMP  Inner        Turtle N TN_Cr~ 3        Crest   3    
##  6   666         2022 GBRMP  Inner        Turtle N TN_Cr~ 3        Crest   8    
##  7   670         2022 GBRMP  Inner        Turtle N TN_Cr~ 3        Crest   12   
##  8   671         2022 GBRMP  Inner        Turtle N TN_Cr~ 3        Crest   13   
##  9   675         2022 GBRMP  Inner        Turtle S TS_Ch~ 5        Crest   3    
## 10   677         2022 GBRMP  Inner        Turtle S TS_Ch~ 5        Crest   5    
## # i 44 more rows
## # i 21 more variables: Acropora <dbl>, Pocillopora <dbl>, Porites <dbl>,
## #   Other <dbl>, total <dbl>, Top_Acr <dbl>, Top_Pocil <dbl>, Top_Por <dbl>,
## #   Top_other <dbl>, Under_Acr <dbl>, Under_Pocil <dbl>, Under_Por <dbl>,
## #   Under_oth <dbl>, Side_Acr <dbl>, Side_Pocil <dbl>, Side_Por <dbl>,
## #   Side_other <dbl>, `Top Total` <dbl>, `Under Total` <dbl>,
## #   `Side Total` <dbl>, notes <chr>
\end{verbatim}

\begin{Shaded}
\begin{Highlighting}[]
\NormalTok{Coral\_settlement\_nGBR }\OtherTok{\textless{}{-}}\NormalTok{ Coral\_settlement\_nGBR }\SpecialCharTok{\%\textgreater{}\%} 
  \FunctionTok{rename}\NormalTok{(}\StringTok{"Shelf\_Position"}\OtherTok{=} \StringTok{\textasciigrave{}}\AttributeTok{Shelf posn}\StringTok{\textasciigrave{}}\NormalTok{,}
         \StringTok{"Year\_Collected"}\OtherTok{=} \StringTok{\textasciigrave{}}\AttributeTok{Year (Jan)}\StringTok{\textasciigrave{}}\NormalTok{,}
         \StringTok{"Site\_\#"}\OtherTok{=} \StringTok{\textasciigrave{}}\AttributeTok{Site \#}\StringTok{\textasciigrave{}}\NormalTok{,}
         \StringTok{"Acroporidae"} \OtherTok{=}\NormalTok{Acropora,}
         \StringTok{"Pocilloporidae"} \OtherTok{=}\NormalTok{ Pocillopora,}
         \StringTok{"Poritid"} \OtherTok{=}\NormalTok{ Porites) }\SpecialCharTok{|\textgreater{}} 
  \FunctionTok{mutate}\NormalTok{(}\AttributeTok{Year\_Collected =} \FunctionTok{factor}\NormalTok{(Year\_Collected))}

\NormalTok{Coral\_settlement\_nGBR }\OtherTok{\textless{}{-}}\NormalTok{ Coral\_settlement\_nGBR }\SpecialCharTok{|\textgreater{}}
  \FunctionTok{filter}\NormalTok{(Shelf\_Position }\SpecialCharTok{==} \StringTok{"Outer"} \SpecialCharTok{|}\NormalTok{ (Shelf\_Position }\SpecialCharTok{\%in\%} \FunctionTok{c}\NormalTok{(}\StringTok{"Inner"}\NormalTok{, }\StringTok{"Mid"}\NormalTok{) }\SpecialCharTok{\&}\NormalTok{ Tile }\SpecialCharTok{\%in\%} \FunctionTok{as.character}\NormalTok{(}\DecValTok{1}\SpecialCharTok{:}\DecValTok{10}\NormalTok{)))}

\NormalTok{Coral\_settlement\_nGBR}\SpecialCharTok{$}\NormalTok{Year\_Collected }\OtherTok{\textless{}{-}} \FunctionTok{factor}\NormalTok{(Coral\_settlement\_nGBR}\SpecialCharTok{$}\NormalTok{Year\_Collected)}
\end{Highlighting}
\end{Shaded}

\begin{Shaded}
\begin{Highlighting}[]
\NormalTok{total\_recruits}\SpecialCharTok{$}\NormalTok{Year\_Collected }\OtherTok{\textless{}{-}} \FunctionTok{as.factor}\NormalTok{(total\_recruits}\SpecialCharTok{$}\NormalTok{Year\_Collected)}
\end{Highlighting}
\end{Shaded}

\subsection{Exploratory analysis}\label{exploratory-analysis}

\begin{Shaded}
\begin{Highlighting}[]
\NormalTok{Fig}\FloatTok{.1}\OtherTok{\textless{}{-}}\NormalTok{  Coral\_settlement\_nGBR }\SpecialCharTok{|\textgreater{}}  
  \FunctionTok{ggplot}\NormalTok{(}\FunctionTok{aes}\NormalTok{(}\AttributeTok{x=}\NormalTok{Year\_Collected, }\AttributeTok{y=}\NormalTok{total, }\AttributeTok{colour =}\NormalTok{ Shelf\_Position))}\SpecialCharTok{+}
  \FunctionTok{geom\_boxplot}\NormalTok{(}\AttributeTok{position =} \StringTok{"dodge"}\NormalTok{)}\SpecialCharTok{+}
  \FunctionTok{theme\_classic}\NormalTok{()}\SpecialCharTok{+}
  \FunctionTok{labs}\NormalTok{(}
  \AttributeTok{y =} \FunctionTok{expression}\NormalTok{(}\StringTok{"Recruits ("} \SpecialCharTok{*}\NormalTok{ tile}\SpecialCharTok{\^{}}\NormalTok{\{}\SpecialCharTok{{-}}\DecValTok{1}\NormalTok{\}}\SpecialCharTok{*}\StringTok{")"}\NormalTok{ ),}
  \AttributeTok{x =} \StringTok{"Year"}\NormalTok{,}
  \AttributeTok{colour =} \StringTok{"Shelf position"}\NormalTok{)}\SpecialCharTok{+}
  \FunctionTok{scale\_colour\_manual}\NormalTok{(}\AttributeTok{values =} \FunctionTok{c}\NormalTok{(}\StringTok{"\#99CCFF"}\NormalTok{, }\StringTok{"\#3399FF"}\NormalTok{, }\StringTok{"\#003366"}\NormalTok{))}\SpecialCharTok{+}
  \FunctionTok{facet\_wrap}\NormalTok{(}\SpecialCharTok{\textasciitilde{}}\NormalTok{Shelf\_Position)}
\end{Highlighting}
\end{Shaded}

Fig 1

\pandocbounded{\includegraphics[keepaspectratio]{Data_Analysis_Thesis_files/figure-latex/Detect Outliers-1.pdf}}

\begin{verbatim}
## # A tibble: 92 x 8
## # Groups:   Family, Shelf_Position, Year_Collected [48]
##    Family         Shelf_Position Year_Collected Reef   average    se   max   min
##    <fct>          <chr>          <fct>          <chr>    <dbl> <dbl> <dbl> <dbl>
##  1 Acroporidae    Inner          2024           Turtl~  136.   24.2    429     5
##  2 Acroporidae    Inner          2024           Turtl~   64.9   8.07   165    10
##  3 Acroporidae    Inner          2022           Turtl~   56.2  11.2    200     0
##  4 Acroporidae    Mid            2024           Lizard   42.3   9.39   223     0
##  5 Acroporidae    Inner          2025           Turtl~   37.1   8.53   243     0
##  6 Acroporidae    Inner          2022           Turtl~   34.3   7.27   139     0
##  7 Acroporidae    Mid            2024           MacGi~   26.6   4.20    90     1
##  8 Acroporidae    Inner          2025           Turtl~   24.6   4.08    69     0
##  9 Pocilloporidae Outer          2024           Hicks    14.2   2.27    45     0
## 10 Acroporidae    Inner          2023           Turtl~   12.4   2.13    43     0
## 11 Acroporidae    Mid            2023           MacGi~   11.7   2.17    37     0
## 12 Acroporidae    Outer          2024           Hicks    11.5   2.28    39     0
## 13 Other          Inner          2024           Turtl~   11.5   1.71    32     1
## 14 Acroporidae    Mid            2025           Lizard   11.4   2.81    62     0
## 15 Acroporidae    Mid            2025           MacGi~    9.86  2.21    47     0
## 16 Pocilloporidae Outer          2025           Hicks     9.62  1.68    37     0
## 17 Pocilloporidae Mid            2022           MacGi~    9.21  2.21    48     0
## 18 Pocilloporidae Mid            2023           MacGi~    8.72  1.45    37     0
## 19 Acroporidae    Mid            2023           Lizard    8.7   1.80    44     0
## 20 Acroporidae    Inner          2023           Turtl~    8.47  1.46    35     0
## # i 72 more rows
\end{verbatim}

\section{Models}\label{models}

\subsection{Question one. is the total number of recruits different
across the
shelf?}\label{question-one.-is-the-total-number-of-recruits-different-across-the-shelf}

\subsubsection{Use GLMM for total
recruits}\label{use-glmm-for-total-recruits}

Run generalised linear models with the Poisson, normal binomail and the
zero inflated normal bionomial. Total is the dependent variable (the
count of coral recruits at a given sampling unit) Shelf position is a
fixed effect predictor, it is used to test if the average coral
settlement differs between shelf positions Random effects: a nested
random effect, modelling the variations between reefs, and variation
between sites nested within reefs

\textbf{\emph{Check the levels of grouping variables}}

\begin{verbatim}
## # A tibble: 18 x 3
##    Reef         Site               n
##    <chr>        <chr>          <int>
##  1 Day          Day 1             38
##  2 Day          Day 2             38
##  3 Day          Day 3             38
##  4 Hicks        Hicks 1           38
##  5 Hicks        Hicks 2           37
##  6 Hicks        Hicks 3           40
##  7 Lizard       Lag2-Trimodal     30
##  8 Lizard       TurtleB-CooksP    30
##  9 Lizard       Vicky-HS          30
## 10 MacGillivray Mac 1             37
## 11 MacGillivray Mac 2             38
## 12 MacGillivray Mac 3             36
## 13 Turtle N     TN_Back           40
## 14 Turtle N     TN_Crest          40
## 15 Turtle N     TN_coral          40
## 16 Turtle S     TS_Back           40
## 17 Turtle S     TS_Channel        40
## 18 Turtle S     TS_Crest          40
\end{verbatim}

All reefs have multiple sites, so nesting should work.

\begin{Shaded}
\begin{Highlighting}[]
\NormalTok{pois\_total\_yr }\OtherTok{\textless{}{-}} \FunctionTok{glmer}\NormalTok{(total }\SpecialCharTok{\textasciitilde{}}\NormalTok{ Shelf\_Position }\SpecialCharTok{+}\NormalTok{ Year\_Collected }\SpecialCharTok{+}\NormalTok{ (}\DecValTok{1} \SpecialCharTok{|}\NormalTok{ Reef}\SpecialCharTok{/}\NormalTok{Site), }
      \AttributeTok{family =} \StringTok{"poisson"}\NormalTok{, }\AttributeTok{data =}\NormalTok{ Coral\_settlement\_nGBR)}

\NormalTok{nbin\_total\_yr }\OtherTok{\textless{}{-}} \FunctionTok{glmmTMB}\NormalTok{(total }\SpecialCharTok{\textasciitilde{}}\NormalTok{ Shelf\_Position }\SpecialCharTok{+}\NormalTok{ Year\_Collected }\SpecialCharTok{+}\NormalTok{ (}\DecValTok{1} \SpecialCharTok{|}\NormalTok{ Reef}\SpecialCharTok{/}\NormalTok{Site), }
        \AttributeTok{family =} \StringTok{"nbinom2"}\NormalTok{, }\AttributeTok{data =}\NormalTok{ Coral\_settlement\_nGBR)}

\NormalTok{zbin\_total\_yr }\OtherTok{\textless{}{-}} \FunctionTok{glmmTMB}\NormalTok{(total }\SpecialCharTok{\textasciitilde{}}\NormalTok{ Shelf\_Position }\SpecialCharTok{+}\NormalTok{ Year\_Collected }\SpecialCharTok{+}\NormalTok{ (}\DecValTok{1} \SpecialCharTok{|}\NormalTok{ Reef}\SpecialCharTok{/}\NormalTok{Site),}
                      \AttributeTok{ziformula =} \SpecialCharTok{\textasciitilde{}}\DecValTok{1}\NormalTok{,}
        \AttributeTok{family =} \StringTok{"nbinom2"}\NormalTok{, }\AttributeTok{data =}\NormalTok{ Coral\_settlement\_nGBR)}

\FunctionTok{AIC}\NormalTok{(pois\_total\_yr, nbin\_total\_yr, zbin\_total\_yr) }\CommentTok{\#nbin has the lowest AIC}
\end{Highlighting}
\end{Shaded}

\begin{verbatim}
##               df       AIC
## pois_total_yr  8 14688.641
## nbin_total_yr  9  5347.038
## zbin_total_yr 10  5348.632
\end{verbatim}

normal binomial model has the lowest AIC so I will get the summary of
that model.

\begin{Shaded}
\begin{Highlighting}[]
\NormalTok{Coral\_settlement\_nGBR }\SpecialCharTok{|\textgreater{}} 
\FunctionTok{group\_by}\NormalTok{(Reef, Site) }\SpecialCharTok{|\textgreater{}} 
  \FunctionTok{summarise}\NormalTok{(}\AttributeTok{count =} \FunctionTok{n}\NormalTok{())}
\end{Highlighting}
\end{Shaded}

\begin{verbatim}
## # A tibble: 18 x 3
## # Groups:   Reef [6]
##    Reef         Site           count
##    <chr>        <chr>          <int>
##  1 Day          Day 1             38
##  2 Day          Day 2             38
##  3 Day          Day 3             38
##  4 Hicks        Hicks 1           38
##  5 Hicks        Hicks 2           37
##  6 Hicks        Hicks 3           40
##  7 Lizard       Lag2-Trimodal     30
##  8 Lizard       TurtleB-CooksP    30
##  9 Lizard       Vicky-HS          30
## 10 MacGillivray Mac 1             37
## 11 MacGillivray Mac 2             38
## 12 MacGillivray Mac 3             36
## 13 Turtle N     TN_Back           40
## 14 Turtle N     TN_Crest          40
## 15 Turtle N     TN_coral          40
## 16 Turtle S     TS_Back           40
## 17 Turtle S     TS_Channel        40
## 18 Turtle S     TS_Crest          40
\end{verbatim}

\begin{verbatim}
##  Family: nbinom2  ( log )
## Formula:          total ~ Shelf_Position + Year_Collected + (1 | Reef/Site)
## Data: Coral_settlement_nGBR
## 
##       AIC       BIC    logLik -2*log(L)  df.resid 
##    5347.0    5387.6   -2664.5    5329.0       661 
## 
## Random effects:
## 
## Conditional model:
##  Groups    Name        Variance  Std.Dev. 
##  Site:Reef (Intercept) 3.436e-01 0.5861576
##  Reef      (Intercept) 6.318e-08 0.0002514
## Number of obs: 670, groups:  Site:Reef, 18; Reef, 6
## 
## Dispersion parameter for nbinom2 family (): 1.28 
## 
## Conditional model:
##                     Estimate Std. Error z value Pr(>|z|)    
## (Intercept)          3.60820    0.25551  14.121  < 2e-16 ***
## Shelf_PositionMid   -0.88736    0.35037  -2.533  0.01132 *  
## Shelf_PositionOuter -1.64001    0.34955  -4.692 2.71e-06 ***
## Year_Collected2023  -0.33985    0.11231  -3.026  0.00248 ** 
## Year_Collected2024   0.85180    0.10809   7.881 3.25e-15 ***
## Year_Collected2025  -0.06039    0.10993  -0.549  0.58280    
## ---
## Signif. codes:  0 '***' 0.001 '**' 0.01 '*' 0.05 '.' 0.1 ' ' 1
\end{verbatim}

2023 and 2024 are significantly different to 2022. 2025 is not
significantly different to 2022. Outer shelf is significantly different
to inner shelf (p\textless0.01), Mid shelf is significantly different to
inner shelf reefs (p\textless0.05). Intercept is also significant: not
quite sure what that means in a biological framework. When making a
figure should I then be using summary statistics or raw count data?
Because the large number of zeros will make it difficult to do a boxplot
of the raw data.

\begin{Shaded}
\begin{Highlighting}[]
\FunctionTok{library}\NormalTok{(parameters)}
\FunctionTok{print}\NormalTok{(}\FunctionTok{model\_parameters}\NormalTok{(nbin\_total\_yr,}
                       \AttributeTok{exponentiate =} \ConstantTok{TRUE} \CommentTok{\#required for log or logit function}
\NormalTok{                       ))}
\end{Highlighting}
\end{Shaded}

\begin{verbatim}
## # Fixed Effects
## 
## Parameter              |   IRR |   SE |         95% CI |     z |      p
## -----------------------------------------------------------------------
## (Intercept)            | 36.90 | 9.43 | [22.36, 60.89] | 14.12 | < .001
## Shelf Position [Mid]   |  0.41 | 0.14 | [ 0.21,  0.82] | -2.53 | 0.011 
## Shelf Position [Outer] |  0.19 | 0.07 | [ 0.10,  0.38] | -4.69 | < .001
## Year Collected [2023]  |  0.71 | 0.08 | [ 0.57,  0.89] | -3.03 | 0.002 
## Year Collected [2024]  |  2.34 | 0.25 | [ 1.90,  2.90] |  7.88 | < .001
## Year Collected [2025]  |  0.94 | 0.10 | [ 0.76,  1.17] | -0.55 | 0.583 
## 
## # Dispersion
## 
## Parameter   | Coefficient |       95% CI
## ----------------------------------------
## (Intercept) |        1.28 | [1.14, 1.44]
## 
## # Random Effects Variances
## 
## Parameter                 | Coefficient |           95% CI
## ----------------------------------------------------------
## SD (Intercept: Site:Reef) |        0.59 | [0.41,     0.83]
## SD (Intercept: Reef)      |    2.51e-04 | [0.00,      Inf]
\end{verbatim}

\begin{Shaded}
\begin{Highlighting}[]
\CommentTok{\# Specify \textquotesingle{}exponentiate = TRUE\textquotesingle{} to interpret coefficients as ratios}
\NormalTok{nbin\_total\_yr.results }\OtherTok{\textless{}{-}}\NormalTok{ parameters}\SpecialCharTok{::}\FunctionTok{model\_parameters}\NormalTok{(nbin\_total\_yr, }\AttributeTok{exponentiate =} \ConstantTok{TRUE}\NormalTok{) }
\end{Highlighting}
\end{Shaded}

\subsubsection{Simulations of the model}\label{simulations-of-the-model}

Simulating model -\textgreater{} need to check the reason people do
this\ldots{} I think it is to see if the observed values are explained
by the model or not.

\begin{Shaded}
\begin{Highlighting}[]
\NormalTok{sim1 }\OtherTok{\textless{}{-}} \FunctionTok{simulateResiduals}\NormalTok{(nbin\_total\_yr)}
\FunctionTok{testResiduals}\NormalTok{(sim1)}
\end{Highlighting}
\end{Shaded}

\pandocbounded{\includegraphics[keepaspectratio]{Data_Analysis_Thesis_files/figure-latex/unnamed-chunk-14-1.pdf}}

\begin{verbatim}
## $uniformity
## 
##  Asymptotic one-sample Kolmogorov-Smirnov test
## 
## data:  simulationOutput$scaledResiduals
## D = 0.02742, p-value = 0.6949
## alternative hypothesis: two-sided
## 
## 
## $dispersion
## 
##  DHARMa nonparametric dispersion test via sd of residuals fitted vs.
##  simulated
## 
## data:  simulationOutput
## dispersion = 0.66066, p-value = 0.728
## alternative hypothesis: two.sided
## 
## 
## $outliers
## 
##  DHARMa outlier test based on exact binomial test with approximate
##  expectations
## 
## data:  simulationOutput
## outliers at both margin(s) = 3, observations = 670, p-value = 0.5079
## alternative hypothesis: true probability of success is not equal to 0.007968127
## 95 percent confidence interval:
##  0.0009243448 0.0130292688
## sample estimates:
## frequency of outliers (expected: 0.00796812749003984 ) 
##                                            0.004477612
\end{verbatim}

\textbf{Key take aways:} - The uniformity indicates the scaled residuals
are uniformly distributed. The p value \textgreater{} 0.05 suggests
there is no evidence the residuals deviate from uniformity. - Dispersion
is close to 1 and p value is \textgreater0.05, suggesting that there is
no overdispersion or under dispersion. - Outliers: suggesting that the
number of outliers is not significantly different to what would be
expected.

Because the simulation suggests not overdispersed, that means I can
either do Wald Z or chi squared test (Bolker, Brooks, Clark, Geange,
Poulsen, Stevens, White; 2008
\url{http://dx.doi.org/10.1016/j.tree.2008.10.008})

\section{Q2 is the composition of the recruits different across the
shelf?}\label{q2-is-the-composition-of-the-recruits-different-across-the-shelf}

PERMANOVA - variation within a group can be calculated directly from a
distance matrix.

Notes from Victor email: Regarding the PERMANOVA: you start with a
community matrix (sites x coral taxa). A PERMANOVA tests whether
community composition differs among groups, using a distance matrix
derived from your multivariate raw data. The inputs are a dissimilarity
(distance) matrix, and a grouping factor (i.e., a categorical variable
you want to test, like ``reef'', ``shelf.position'', or ``year''). In
your case, I'd suggest calculating a Bray-Curtis distance matrix and
then running a PERMANOVA with the vegan package. Here's a quick example
you can run in R: library(vegan)

\section{Example data: 6 sites x 4 species (abundance
table)}\label{example-data-6-sites-x-4-species-abundance-table}

abund \textless- data.frame(

sp1 = c(5, 3, 0, 2, 7, 1),

sp2 = c(0, 1, 2, 3, 0, 4),

sp3 = c(8, 6, 5, 9, 7, 6),

sp4 = c(1, 0, 2, 1, 0, 3))

rownames(abund) \textless- paste0(``site'', 1:6)

\section{Metadata with two grouping
factors}\label{metadata-with-two-grouping-factors}

metadata \textless- data.frame(

site = rownames(abund),

region = rep(c(``north'', ``south''), each = 5), \# Factor 1

shelf = rep(c(``inshore'', ``mid''), times = 5) \# Factor 2

dist.mat \textless- vegdist(abund, method = ``bray'') \# Step 1:
Calculate dissimilarity matrix (Bray-Curtis)

adonis2(dist.mat \textasciitilde{} region * shelf, data = metadata,
permutations = 999) \# Step 2: run PERMANOVA

\subsection{NMDS}\label{nmds}

Asking the question: how does the coral recruitment composition differ
in vegetated tiles\ldots{}

\begin{Shaded}
\begin{Highlighting}[]
\FunctionTok{library}\NormalTok{(vegan)}

\CommentTok{\#Create a site{-}by{-}species matrix which includes year and reef as the "site"}

\CommentTok{\# Step 1: Create Site column and keep only species data}
\NormalTok{site\_by\_species }\OtherTok{\textless{}{-}}\NormalTok{ Coral\_settlement\_nGBR }\SpecialCharTok{\%\textgreater{}\%} 
  \FunctionTok{select}\NormalTok{(Year\_Collected, Reef, Acroporidae, Pocilloporidae, Poritid, Other) }\SpecialCharTok{\%\textgreater{}\%} 
  \FunctionTok{unite}\NormalTok{(}\StringTok{"Site"}\NormalTok{, Year\_Collected, Reef, }\AttributeTok{sep =} \StringTok{"\_"}\NormalTok{)}

\CommentTok{\# Step 2: Store species{-}only matrix separately}
\NormalTok{species\_matrix }\OtherTok{\textless{}{-}}\NormalTok{ site\_by\_species }\SpecialCharTok{\%\textgreater{}\%}
  \FunctionTok{select}\NormalTok{(Acroporidae, Pocilloporidae, Poritid, Other)}

\CommentTok{\# Remove rows where all species counts are 0}
\NormalTok{species\_matrix\_filtered }\OtherTok{\textless{}{-}}\NormalTok{ species\_matrix[}\FunctionTok{rowSums}\NormalTok{(species\_matrix) }\SpecialCharTok{\textgreater{}} \DecValTok{0}\NormalTok{, ] }\CommentTok{\#this step removes 43 tiles that have no recruits.}

\CommentTok{\# Step 3: Run NMDS}
\NormalTok{NMDS }\OtherTok{\textless{}{-}} \FunctionTok{metaMDS}\NormalTok{(species\_matrix\_filtered, }\AttributeTok{distance =} \StringTok{"bray"}\NormalTok{, }\AttributeTok{autotransform =} \ConstantTok{FALSE}\NormalTok{, }\AttributeTok{k =} \DecValTok{2}\NormalTok{, }\AttributeTok{trymax =} \DecValTok{100}\NormalTok{)}
\end{Highlighting}
\end{Shaded}

\begin{verbatim}
## Run 0 stress 0.1745302 
## Run 1 stress 0.2001942 
## Run 2 stress 0.197038 
## Run 3 stress 0.19043 
## Run 4 stress 0.1878581 
## Run 5 stress 0.1823981 
## Run 6 stress 0.1973656 
## Run 7 stress 0.1948448 
## Run 8 stress 0.1854282 
## Run 9 stress 0.1939996 
## Run 10 stress 0.192629 
## Run 11 stress 0.1802377 
## Run 12 stress 0.1975926 
## Run 13 stress 0.2011235 
## Run 14 stress 0.186324 
## Run 15 stress 0.1966268 
## Run 16 stress 0.19393 
## Run 17 stress 0.2078842 
## Run 18 stress 0.1867171 
## Run 19 stress 0.1877302 
## Run 20 stress 0.1815839 
## Run 21 stress 0.197914 
## Run 22 stress 0.1867361 
## Run 23 stress 0.2194766 
## Run 24 stress 0.1994805 
## Run 25 stress 0.1925438 
## Run 26 stress 0.1967635 
## Run 27 stress 0.2056334 
## Run 28 stress 0.1835677 
## Run 29 stress 0.1937144 
## Run 30 stress 0.2054904 
## Run 31 stress 0.1838726 
## Run 32 stress 0.1874636 
## Run 33 stress 0.2006946 
## Run 34 stress 0.1925334 
## Run 35 stress 0.1919332 
## Run 36 stress 0.1792134 
## Run 37 stress 0.1797259 
## Run 38 stress 0.1752724 
## Run 39 stress 0.1838911 
## Run 40 stress 0.1835544 
## Run 41 stress 0.1857579 
## Run 42 stress 0.1950891 
## Run 43 stress 0.1823317 
## Run 44 stress 0.1754317 
## Run 45 stress 0.1994145 
## Run 46 stress 0.1920592 
## Run 47 stress 0.192306 
## Run 48 stress 0.1884891 
## Run 49 stress 0.2083178 
## Run 50 stress 0.1955837 
## Run 51 stress 0.1958095 
## Run 52 stress 0.1814009 
## Run 53 stress 0.1987653 
## Run 54 stress 0.1916601 
## Run 55 stress 0.1807283 
## Run 56 stress 0.1924755 
## Run 57 stress 0.1903739 
## Run 58 stress 0.1908093 
## Run 59 stress 0.1952725 
## Run 60 stress 0.1917557 
## Run 61 stress 0.181882 
## Run 62 stress 0.1775983 
## Run 63 stress 0.1956638 
## Run 64 stress 0.1775816 
## Run 65 stress 0.1910875 
## Run 66 stress 0.1861129 
## Run 67 stress 0.2097842 
## Run 68 stress 0.2043553 
## Run 69 stress 0.1941095 
## Run 70 stress 0.1963199 
## Run 71 stress 0.2136376 
## Run 72 stress 0.1838085 
## Run 73 stress 0.1949394 
## Run 74 stress 0.1855656 
## Run 75 stress 0.1796715 
## Run 76 stress 0.1908532 
## Run 77 stress 0.2143244 
## Run 78 stress 0.2030811 
## Run 79 stress 0.1832611 
## Run 80 stress 0.1937234 
## Run 81 stress 0.1821389 
## Run 82 stress 0.1795855 
## Run 83 stress 0.2098308 
## Run 84 stress 0.2027357 
## Run 85 stress 0.1924952 
## Run 86 stress 0.2016489 
## Run 87 stress 0.1788048 
## Run 88 stress 0.2156538 
## Run 89 stress 0.1964815 
## Run 90 stress 0.1753465 
## Run 91 stress 0.1814067 
## Run 92 stress 0.2075672 
## Run 93 stress 0.1814972 
## Run 94 stress 0.1951835 
## Run 95 stress 0.1782501 
## Run 96 stress 0.1833717 
## Run 97 stress 0.2110453 
## Run 98 stress 0.1934076 
## Run 99 stress 0.2037349 
## Run 100 stress 0.1860983 
## *** Best solution was not repeated -- monoMDS stopping criteria:
##      6: no. of iterations >= maxit
##     34: stress ratio > sratmax
##     60: scale factor of the gradient < sfgrmin
\end{verbatim}

\begin{Shaded}
\begin{Highlighting}[]
\CommentTok{\#NMDS\_3 \textless{}{-} metaMDS(species\_matrix\_filtered, distance = "bray", autotransform = FALSE, k = 3, trymax = 100)}


\FunctionTok{print}\NormalTok{(NMDS)}
\end{Highlighting}
\end{Shaded}

\begin{verbatim}
## 
## Call:
## metaMDS(comm = species_matrix_filtered, distance = "bray", k = 2,      trymax = 100, autotransform = FALSE) 
## 
## global Multidimensional Scaling using monoMDS
## 
## Data:     species_matrix_filtered 
## Distance: bray 
## 
## Dimensions: 2 
## Stress:     0.1745302 
## Stress type 1, weak ties
## Best solution was not repeated after 100 tries
## The best solution was from try 0 (metric scaling or null solution)
## Scaling: centring, PC rotation, halfchange scaling 
## Species: expanded scores based on 'species_matrix_filtered'
\end{verbatim}

\begin{Shaded}
\begin{Highlighting}[]
\CommentTok{\#print(NMDS\_3)}

\NormalTok{NMDS}\SpecialCharTok{$}\NormalTok{stress}
\end{Highlighting}
\end{Shaded}

\begin{verbatim}
## [1] 0.1745302
\end{verbatim}

\begin{Shaded}
\begin{Highlighting}[]
\CommentTok{\#NMDS\_3$stress}
\CommentTok{\#the stress score is 0.118}
\end{Highlighting}
\end{Shaded}

Assess the fit of a NMDS by plotting the original dissimilarities (diss)
against the Euclidean ordination distances (dist). Shepard plot

\begin{Shaded}
\begin{Highlighting}[]
\FunctionTok{plot}\NormalTok{(NMDS}\SpecialCharTok{$}\NormalTok{diss, NMDS}\SpecialCharTok{$}\NormalTok{dist)}
\end{Highlighting}
\end{Shaded}

\pandocbounded{\includegraphics[keepaspectratio]{Data_Analysis_Thesis_files/figure-latex/unnamed-chunk-16-1.pdf}}

\begin{Shaded}
\begin{Highlighting}[]
\FunctionTok{stressplot}\NormalTok{(}\AttributeTok{object =}\NormalTok{ NMDS)}
\end{Highlighting}
\end{Shaded}

\pandocbounded{\includegraphics[keepaspectratio]{Data_Analysis_Thesis_files/figure-latex/unnamed-chunk-16-2.pdf}}

NMDS ordinations are often plotted without units or labels on the axes
so that people do not try to interpret them.

\begin{Shaded}
\begin{Highlighting}[]
\FunctionTok{ordiplot}\NormalTok{(NMDS, }\AttributeTok{type =} \StringTok{"n"}\NormalTok{)}
\FunctionTok{orditorp}\NormalTok{(NMDS, }\AttributeTok{display =} \StringTok{"species"}\NormalTok{, }\AttributeTok{col =} \StringTok{"blue"}\NormalTok{, }\AttributeTok{air =} \FloatTok{0.01}\NormalTok{)}
\FunctionTok{orditorp}\NormalTok{(NMDS, }\AttributeTok{display =} \StringTok{"sites"}\NormalTok{, }\AttributeTok{cex =} \FloatTok{0.5}\NormalTok{, }\AttributeTok{air =} \FloatTok{0.01}\NormalTok{)}
\end{Highlighting}
\end{Shaded}

\pandocbounded{\includegraphics[keepaspectratio]{Data_Analysis_Thesis_files/figure-latex/unnamed-chunk-17-1.pdf}}

\subsubsection{PERMANOVA}\label{permanova}

\begin{Shaded}
\begin{Highlighting}[]
\FunctionTok{library}\NormalTok{(vegan)}
\FunctionTok{library}\NormalTok{(dplyr)}
\FunctionTok{library}\NormalTok{(tidyr)}

\CommentTok{\# Step 1: Create site{-}by{-}species matrix and "Site" column}
\NormalTok{site\_by\_species }\OtherTok{\textless{}{-}}\NormalTok{ Coral\_settlement\_nGBR }\SpecialCharTok{\%\textgreater{}\%}
  \FunctionTok{select}\NormalTok{(Year\_Collected, Reef, Shelf\_Position,Site, Tile, Acroporidae, Pocilloporidae, Poritid, Other) }

\CommentTok{\# Add a unique Site ID *after* extracting metadata}
\NormalTok{site\_by\_species }\OtherTok{\textless{}{-}}\NormalTok{ site\_by\_species }\SpecialCharTok{\%\textgreater{}\%}
  \FunctionTok{mutate}\NormalTok{(}\AttributeTok{Site\_ID =} \FunctionTok{paste}\NormalTok{(Year\_Collected, Reef, Site, Tile, }\AttributeTok{sep =} \StringTok{"\_"}\NormalTok{))}

\CommentTok{\# Step 2: Separate metadata and species data}
\NormalTok{metadata }\OtherTok{\textless{}{-}}\NormalTok{ site\_by\_species }\SpecialCharTok{\%\textgreater{}\%}
  \FunctionTok{select}\NormalTok{(Site\_ID, Shelf\_Position, Year\_Collected)}

\NormalTok{species\_matrix }\OtherTok{\textless{}{-}}\NormalTok{ site\_by\_species }\SpecialCharTok{\%\textgreater{}\%}
  \FunctionTok{select}\NormalTok{(Site\_ID, Acroporidae, Pocilloporidae, Poritid, Other)}

\CommentTok{\# Step 3: Set row names and remove species rows with all zeros}
\NormalTok{species\_matrix\_filtered }\OtherTok{\textless{}{-}}\NormalTok{ species\_matrix }\SpecialCharTok{\%\textgreater{}\%}
  \FunctionTok{column\_to\_rownames}\NormalTok{(}\StringTok{"Site\_ID"}\NormalTok{) }\SpecialCharTok{\%\textgreater{}\%}
  \FunctionTok{filter}\NormalTok{(}\FunctionTok{rowSums}\NormalTok{(.) }\SpecialCharTok{\textgreater{}} \DecValTok{0}\NormalTok{)}

\CommentTok{\# Step 4: Match metadata rows to filtered species matrix}
\NormalTok{metadata\_filtered }\OtherTok{\textless{}{-}}\NormalTok{ metadata }\SpecialCharTok{\%\textgreater{}\%}
  \FunctionTok{filter}\NormalTok{(Site\_ID }\SpecialCharTok{\%in\%} \FunctionTok{rownames}\NormalTok{(species\_matrix\_filtered)) }\SpecialCharTok{\%\textgreater{}\%}
  \FunctionTok{column\_to\_rownames}\NormalTok{(}\StringTok{"Site\_ID"}\NormalTok{)}

\CommentTok{\# Step 5: PERMANOVA}
\NormalTok{adonis\_result }\OtherTok{\textless{}{-}}\NormalTok{ vegan}\SpecialCharTok{::}\FunctionTok{adonis2}\NormalTok{(species\_matrix\_filtered }\SpecialCharTok{\textasciitilde{}}\NormalTok{ Shelf\_Position}\SpecialCharTok{+}\NormalTok{ Year\_Collected,}
                                \AttributeTok{data =}\NormalTok{ metadata\_filtered,}
                                \AttributeTok{method =} \StringTok{"bray"}\NormalTok{,}
                                \AttributeTok{by =} \StringTok{"term"}\NormalTok{,}
                                \AttributeTok{permutations =} \DecValTok{999}\NormalTok{)}

\FunctionTok{print}\NormalTok{(adonis\_result)}
\end{Highlighting}
\end{Shaded}

\begin{verbatim}
## Permutation test for adonis under reduced model
## Terms added sequentially (first to last)
## Permutation: free
## Number of permutations: 999
## 
## vegan::adonis2(formula = species_matrix_filtered ~ Shelf_Position + Year_Collected, data = metadata_filtered, permutations = 999, method = "bray", by = "term")
##                 Df SumOfSqs      R2      F Pr(>F)    
## Shelf_Position   2   30.623 0.18122 72.868  0.001 ***
## Year_Collected   3    7.867 0.04656 12.480  0.001 ***
## Residual       621  130.490 0.77222                  
## Total          626  168.981 1.00000                  
## ---
## Signif. codes:  0 '***' 0.001 '**' 0.01 '*' 0.05 '.' 0.1 ' ' 1
\end{verbatim}

The PERMANOVA suggests that shelf position and year do significantly
affect the coral community composition. Shelf position explains 17.4\%
of the variation (strong and significant) Year collected explains 5\% of
the variation (suggesting significant temporal differences) 77.5\% of
the variation in composition of coral recruits remains unexplained.

\section{Redundancy Analysis}\label{redundancy-analysis}

\url{https://pmassicotte.github.io/stats-denmark-2019/07_rda.html\#/a-concrete-example-3}

Suggests using the Hellinger-transform method (Borcard, Gillet, and
Legendre 2011; Legendre and Gallagher 2001) that is suited to species
abundance data,

\begin{itemize}
\tightlist
\item
  gives low weights to variables with low counts and many zeros. Divide
  each value in a data matrix by its row sum, and taking the square root
  of the quotient
\end{itemize}

\begin{Shaded}
\begin{Highlighting}[]
\NormalTok{Recruits }\OtherTok{\textless{}{-}}\NormalTok{ Coral\_settlement\_nGBR }\SpecialCharTok{\%\textgreater{}\%}
  \FunctionTok{select}\NormalTok{(Acroporidae, Pocilloporidae, Poritid, Other)}
\NormalTok{recruit\_hellinger }\OtherTok{\textless{}{-}} \FunctionTok{decostand}\NormalTok{(Recruits, }\AttributeTok{method =} \StringTok{"hellinger"}\NormalTok{)}

\NormalTok{environment\_data }\OtherTok{\textless{}{-}}\NormalTok{ Coral\_settlement\_nGBR }\SpecialCharTok{\%\textgreater{}\%}
  \FunctionTok{select}\NormalTok{(Year\_Collected, Site, Reef, Shelf\_Position) }\CommentTok{\#worried because environment data is nested here.}
\end{Highlighting}
\end{Shaded}

\begin{Shaded}
\begin{Highlighting}[]
\NormalTok{my\_rda }\OtherTok{\textless{}{-}} \FunctionTok{rda}\NormalTok{(recruit\_hellinger}\SpecialCharTok{\textasciitilde{}}\NormalTok{Year\_Collected }\SpecialCharTok{+}\NormalTok{ Shelf\_Position}\SpecialCharTok{+}\NormalTok{ (Reef}\SpecialCharTok{/}\NormalTok{Site), }\AttributeTok{data =}\NormalTok{ environment\_data)}

\FunctionTok{summary}\NormalTok{(my\_rda)}
\end{Highlighting}
\end{Shaded}

\begin{verbatim}
## 
## Call:
## rda(formula = recruit_hellinger ~ Year_Collected + Shelf_Position +      (Reef/Site), data = environment_data) 
## 
## Partitioning of variance:
##               Inertia Proportion
## Total          0.3727     1.0000
## Constrained    0.1686     0.4525
## Unconstrained  0.2040     0.5475
## 
## Eigenvalues, and their contribution to the variance 
## 
## Importance of components:
##                         RDA1    RDA2     RDA3     RDA4     PC1    PC2     PC3
## Eigenvalue            0.1343 0.02557 0.006092 0.002612 0.07728 0.0593 0.04022
## Proportion Explained  0.3605 0.06863 0.016348 0.007010 0.20737 0.1591 0.10794
## Cumulative Proportion 0.3605 0.42911 0.445462 0.452471 0.65984 0.8190 0.92690
##                           PC4
## Eigenvalue            0.02724
## Proportion Explained  0.07310
## Cumulative Proportion 1.00000
## 
## Accumulated constrained eigenvalues
## Importance of components:
##                         RDA1    RDA2     RDA3     RDA4
## Eigenvalue            0.1343 0.02557 0.006092 0.002612
## Proportion Explained  0.7967 0.15167 0.036131 0.015492
## Cumulative Proportion 0.7967 0.94838 0.984508 1.000000
\end{verbatim}

\begin{Shaded}
\begin{Highlighting}[]
\FunctionTok{plot}\NormalTok{(my\_rda)}
\end{Highlighting}
\end{Shaded}

\pandocbounded{\includegraphics[keepaspectratio]{Data_Analysis_Thesis_files/figure-latex/unnamed-chunk-21-1.pdf}}
Explained variance

\begin{Shaded}
\begin{Highlighting}[]
\FunctionTok{RsquareAdj}\NormalTok{(my\_rda)}
\end{Highlighting}
\end{Shaded}

\begin{verbatim}
## $r.squared
## [1] 0.4524714
## 
## $adj.r.squared
## [1] 0.4355984
\end{verbatim}

Tests of RDA results For canonical analysis such as RDA, model
significance relies on permutation tests.

The principle of a permutation test is to generate a reference
distribution of the chosen statistic under the null hypothesis H0 by
randomly permuting appropriate elements of the data a large number of
times and recomputing the statistic each time. Then, one compares the
true value of the statistic to this reference distribution. The p value
is computed as the proportion of the permuted values equal to or larger
than the true (unpermuted) value of the statistic for a one-tailed test
in the upper tail, like the F test used in RDA (Borcard, Gillet, and
Legendre 2011).

H0: Observed results can be produced by random chance.

The null hypothesis is rejected if this p value is equal to or smaller
than the predefined significance level α (ex.: α=0.05).

Test for 1. Global RDA significance 2. Axis significance 3. Terms
(explanatory variables) significance

\subsubsection{Global significance test}\label{global-significance-test}

\begin{Shaded}
\begin{Highlighting}[]
\FunctionTok{anova.cca}\NormalTok{(my\_rda, }\AttributeTok{permutations =} \DecValTok{999}\NormalTok{)}
\end{Highlighting}
\end{Shaded}

\begin{verbatim}
## Permutation test for rda under reduced model
## Permutation: free
## Number of permutations: 999
## 
## Model: rda(formula = recruit_hellinger ~ Year_Collected + Shelf_Position + (Reef/Site), data = environment_data)
##           Df Variance      F Pr(>F)    
## Model     20  0.16862 26.816  0.001 ***
## Residual 649  0.20404                  
## ---
## Signif. codes:  0 '***' 0.001 '**' 0.01 '*' 0.05 '.' 0.1 ' ' 1
\end{verbatim}

Since p = 0.001, reject H0, meaning that the RDA model is significant
and not produced by random chance

\subsubsection{Axis significance}\label{axis-significance}

\begin{Shaded}
\begin{Highlighting}[]
\FunctionTok{anova.cca}\NormalTok{(my\_rda, }\AttributeTok{by =} \StringTok{"axis"}\NormalTok{)}
\end{Highlighting}
\end{Shaded}

\begin{verbatim}
## Permutation test for rda under reduced model
## Forward tests for axes
## Permutation: free
## Number of permutations: 999
## 
## Model: rda(formula = recruit_hellinger ~ Year_Collected + Shelf_Position + (Reef/Site), data = environment_data)
##           Df Variance        F Pr(>F)    
## RDA1       1 0.134339 427.2946  0.001 ***
## RDA2       1 0.025574  83.4755  0.001 ***
## RDA3       1 0.006092  19.9155  0.001 ***
## RDA4       1 0.002612   8.5522  0.001 ***
## Residual 665 0.204042                    
## ---
## Signif. codes:  0 '***' 0.001 '**' 0.01 '*' 0.05 '.' 0.1 ' ' 1
\end{verbatim}

\subsubsection{Terms significance}\label{terms-significance}

\begin{Shaded}
\begin{Highlighting}[]
\FunctionTok{anova.cca}\NormalTok{(my\_rda, }\AttributeTok{by =} \StringTok{"terms"}\NormalTok{)}
\end{Highlighting}
\end{Shaded}

\begin{verbatim}
## Permutation test for rda under reduced model
## Terms added sequentially (first to last)
## Permutation: free
## Number of permutations: 999
## 
## Model: rda(formula = recruit_hellinger ~ Year_Collected + Shelf_Position + (Reef/Site), data = environment_data)
##                 Df Variance        F Pr(>F)    
## Year_Collected   3 0.010548  11.1833  0.001 ***
## Shelf_Position   2 0.111391 177.1510  0.001 ***
## Reef             3 0.011063  11.7298  0.001 ***
## Reef:Site       12 0.035616   9.4404  0.001 ***
## Residual       649 0.204042                    
## ---
## Signif. codes:  0 '***' 0.001 '**' 0.01 '*' 0.05 '.' 0.1 ' ' 1
\end{verbatim}

\subsubsection{Check for
multicollinearity}\label{check-for-multicollinearity}

Use the vif.cca function

\begin{Shaded}
\begin{Highlighting}[]
\FunctionTok{sqrt}\NormalTok{(}\FunctionTok{vif.cca}\NormalTok{(my\_rda))}
\end{Highlighting}
\end{Shaded}

\begin{verbatim}
##                  Year_Collected2023                  Year_Collected2024 
##                            1.312300                            1.314602 
##                  Year_Collected2025                   Shelf_PositionMid 
##                            1.313481                            2.725301 
##                 Shelf_PositionOuter                           ReefHicks 
##                            2.781270                            2.211065 
##                          ReefLizard                    ReefMacGillivray 
##                            2.190414                                  NA 
##                        ReefTurtle N                        ReefTurtle S 
##                            2.219318                                  NA 
##                   ReefDay:SiteDay 2                 ReefHicks:SiteDay 2 
##                            1.373901                                  NA 
##                ReefLizard:SiteDay 2          ReefMacGillivray:SiteDay 2 
##                                  NA                                  NA 
##              ReefTurtle N:SiteDay 2              ReefTurtle S:SiteDay 2 
##                                  NA                                  NA 
##                   ReefDay:SiteDay 3                 ReefHicks:SiteDay 3 
##                            1.373643                                  NA 
##                ReefLizard:SiteDay 3          ReefMacGillivray:SiteDay 3 
##                                  NA                                  NA 
##              ReefTurtle N:SiteDay 3              ReefTurtle S:SiteDay 3 
##                                  NA                                  NA 
##                 ReefDay:SiteHicks 1               ReefHicks:SiteHicks 1 
##                                  NA                            1.356302 
##              ReefLizard:SiteHicks 1        ReefMacGillivray:SiteHicks 1 
##                                  NA                                  NA 
##            ReefTurtle N:SiteHicks 1            ReefTurtle S:SiteHicks 1 
##                                  NA                                  NA 
##                 ReefDay:SiteHicks 2               ReefHicks:SiteHicks 2 
##                                  NA                            1.348631 
##              ReefLizard:SiteHicks 2        ReefMacGillivray:SiteHicks 2 
##                                  NA                                  NA 
##            ReefTurtle N:SiteHicks 2            ReefTurtle S:SiteHicks 2 
##                                  NA                                  NA 
##                 ReefDay:SiteHicks 3               ReefHicks:SiteHicks 3 
##                                  NA                                  NA 
##              ReefLizard:SiteHicks 3        ReefMacGillivray:SiteHicks 3 
##                                  NA                                  NA 
##            ReefTurtle N:SiteHicks 3            ReefTurtle S:SiteHicks 3 
##                                  NA                                  NA 
##           ReefDay:SiteLag2-Trimodal         ReefHicks:SiteLag2-Trimodal 
##                                  NA                                  NA 
##        ReefLizard:SiteLag2-Trimodal  ReefMacGillivray:SiteLag2-Trimodal 
##                            1.382189                                  NA 
##      ReefTurtle N:SiteLag2-Trimodal      ReefTurtle S:SiteLag2-Trimodal 
##                                  NA                                  NA 
##                   ReefDay:SiteMac 1                 ReefHicks:SiteMac 1 
##                                  NA                                  NA 
##                ReefLizard:SiteMac 1          ReefMacGillivray:SiteMac 1 
##                                  NA                            1.384291 
##              ReefTurtle N:SiteMac 1              ReefTurtle S:SiteMac 1 
##                                  NA                                  NA 
##                   ReefDay:SiteMac 2                 ReefHicks:SiteMac 2 
##                                  NA                                  NA 
##                ReefLizard:SiteMac 2          ReefMacGillivray:SiteMac 2 
##                                  NA                            1.392636 
##              ReefTurtle N:SiteMac 2              ReefTurtle S:SiteMac 2 
##                                  NA                                  NA 
##                   ReefDay:SiteMac 3                 ReefHicks:SiteMac 3 
##                                  NA                                  NA 
##                ReefLizard:SiteMac 3          ReefMacGillivray:SiteMac 3 
##                                  NA                                  NA 
##              ReefTurtle N:SiteMac 3              ReefTurtle S:SiteMac 3 
##                                  NA                                  NA 
##                 ReefDay:SiteTN_Back               ReefHicks:SiteTN_Back 
##                                  NA                                  NA 
##              ReefLizard:SiteTN_Back        ReefMacGillivray:SiteTN_Back 
##                                  NA                                  NA 
##            ReefTurtle N:SiteTN_Back            ReefTurtle S:SiteTN_Back 
##                            1.371349                                  NA 
##                ReefDay:SiteTN_coral              ReefHicks:SiteTN_coral 
##                                  NA                                  NA 
##             ReefLizard:SiteTN_coral       ReefMacGillivray:SiteTN_coral 
##                                  NA                                  NA 
##           ReefTurtle N:SiteTN_coral           ReefTurtle S:SiteTN_coral 
##                            1.371349                                  NA 
##                ReefDay:SiteTN_Crest              ReefHicks:SiteTN_Crest 
##                                  NA                                  NA 
##             ReefLizard:SiteTN_Crest       ReefMacGillivray:SiteTN_Crest 
##                                  NA                                  NA 
##           ReefTurtle N:SiteTN_Crest           ReefTurtle S:SiteTN_Crest 
##                                  NA                                  NA 
##                 ReefDay:SiteTS_Back               ReefHicks:SiteTS_Back 
##                                  NA                                  NA 
##              ReefLizard:SiteTS_Back        ReefMacGillivray:SiteTS_Back 
##                                  NA                                  NA 
##            ReefTurtle N:SiteTS_Back            ReefTurtle S:SiteTS_Back 
##                                  NA                            1.371349 
##              ReefDay:SiteTS_Channel            ReefHicks:SiteTS_Channel 
##                                  NA                                  NA 
##           ReefLizard:SiteTS_Channel     ReefMacGillivray:SiteTS_Channel 
##                                  NA                                  NA 
##         ReefTurtle N:SiteTS_Channel         ReefTurtle S:SiteTS_Channel 
##                                  NA                            1.371349 
##                ReefDay:SiteTS_Crest              ReefHicks:SiteTS_Crest 
##                                  NA                                  NA 
##             ReefLizard:SiteTS_Crest       ReefMacGillivray:SiteTS_Crest 
##                                  NA                                  NA 
##           ReefTurtle N:SiteTS_Crest           ReefTurtle S:SiteTS_Crest 
##                                  NA                                  NA 
##          ReefDay:SiteTurtleB-CooksP        ReefHicks:SiteTurtleB-CooksP 
##                                  NA                                  NA 
##       ReefLizard:SiteTurtleB-CooksP ReefMacGillivray:SiteTurtleB-CooksP 
##                            1.382189                                  NA 
##     ReefTurtle N:SiteTurtleB-CooksP     ReefTurtle S:SiteTurtleB-CooksP 
##                                  NA                                  NA 
##                ReefDay:SiteVicky-HS              ReefHicks:SiteVicky-HS 
##                                  NA                                  NA 
##             ReefLizard:SiteVicky-HS       ReefMacGillivray:SiteVicky-HS 
##                                  NA                                  NA 
##           ReefTurtle N:SiteVicky-HS           ReefTurtle S:SiteVicky-HS 
##                                  NA                                  NA
\end{verbatim}

Not sure what the go is with some being NA. ``As a rule of thumb, if
sqrt vif \textgreater2, collinarity is considered high and some of the
explanatory cariables are redundant.

I will remove reef and site.

\begin{Shaded}
\begin{Highlighting}[]
\NormalTok{my\_rda2 }\OtherTok{\textless{}{-}} \FunctionTok{rda}\NormalTok{(recruit\_hellinger}\SpecialCharTok{\textasciitilde{}}\NormalTok{Year\_Collected }\SpecialCharTok{+}\NormalTok{ Shelf\_Position, }\AttributeTok{data =}\NormalTok{ environment\_data)}

\FunctionTok{summary}\NormalTok{(my\_rda2)}
\end{Highlighting}
\end{Shaded}

\begin{verbatim}
## 
## Call:
## rda(formula = recruit_hellinger ~ Year_Collected + Shelf_Position,      data = environment_data) 
## 
## Partitioning of variance:
##               Inertia Proportion
## Total          0.3727     1.0000
## Constrained    0.1219     0.3272
## Unconstrained  0.2507     0.6728
## 
## Eigenvalues, and their contribution to the variance 
## 
## Importance of components:
##                         RDA1     RDA2     RDA3      RDA4     PC1     PC2
## Eigenvalue            0.1109 0.007936 0.002847 0.0002622 0.09369 0.07859
## Proportion Explained  0.2976 0.021297 0.007639 0.0007036 0.25142 0.21088
## Cumulative Proportion 0.2976 0.318869 0.326508 0.3272111 0.57863 0.78952
##                           PC3     PC4
## Eigenvalue            0.04994 0.02850
## Proportion Explained  0.13401 0.07647
## Cumulative Proportion 0.92353 1.00000
## 
## Accumulated constrained eigenvalues
## Importance of components:
##                         RDA1     RDA2     RDA3      RDA4
## Eigenvalue            0.1109 0.007936 0.002847 0.0002622
## Proportion Explained  0.9094 0.065085 0.023345 0.0021503
## Cumulative Proportion 0.9094 0.974505 0.997850 1.0000000
\end{verbatim}

\begin{Shaded}
\begin{Highlighting}[]
\FunctionTok{plot}\NormalTok{(my\_rda2)}
\end{Highlighting}
\end{Shaded}

\pandocbounded{\includegraphics[keepaspectratio]{Data_Analysis_Thesis_files/figure-latex/unnamed-chunk-28-1.pdf}}

\begin{Shaded}
\begin{Highlighting}[]
\FunctionTok{RsquareAdj}\NormalTok{(my\_rda2)}
\end{Highlighting}
\end{Shaded}

\begin{verbatim}
## $r.squared
## [1] 0.3272111
## 
## $adj.r.squared
## [1] 0.3221449
\end{verbatim}

\begin{Shaded}
\begin{Highlighting}[]
\FunctionTok{anova.cca}\NormalTok{(my\_rda2, }\AttributeTok{permutations =} \DecValTok{999}\NormalTok{)}
\end{Highlighting}
\end{Shaded}

\begin{verbatim}
## Permutation test for rda under reduced model
## Permutation: free
## Number of permutations: 999
## 
## Model: rda(formula = recruit_hellinger ~ Year_Collected + Shelf_Position, data = environment_data)
##           Df Variance      F Pr(>F)    
## Model      5  0.12194 64.587  0.001 ***
## Residual 664  0.25072                  
## ---
## Signif. codes:  0 '***' 0.001 '**' 0.01 '*' 0.05 '.' 0.1 ' ' 1
\end{verbatim}

\begin{Shaded}
\begin{Highlighting}[]
\FunctionTok{anova.cca}\NormalTok{(my\_rda2, }\AttributeTok{by =} \StringTok{"axis"}\NormalTok{)}
\end{Highlighting}
\end{Shaded}

\begin{verbatim}
## Permutation test for rda under reduced model
## Forward tests for axes
## Permutation: free
## Number of permutations: 999
## 
## Model: rda(formula = recruit_hellinger ~ Year_Collected + Shelf_Position, data = environment_data)
##           Df Variance        F Pr(>F)    
## RDA1       1 0.110893 293.6849  0.001 ***
## RDA2       1 0.007936  21.0817  0.001 ***
## RDA3       1 0.002847   7.5730  0.001 ***
## RDA4       1 0.000262   0.6986  0.589    
## Residual 665 0.250722                    
## ---
## Signif. codes:  0 '***' 0.001 '**' 0.01 '*' 0.05 '.' 0.1 ' ' 1
\end{verbatim}

\begin{Shaded}
\begin{Highlighting}[]
\FunctionTok{anova.cca}\NormalTok{(my\_rda2, }\AttributeTok{by =} \StringTok{"terms"}\NormalTok{)}
\end{Highlighting}
\end{Shaded}

\begin{verbatim}
## Permutation test for rda under reduced model
## Terms added sequentially (first to last)
## Permutation: free
## Number of permutations: 999
## 
## Model: rda(formula = recruit_hellinger ~ Year_Collected + Shelf_Position, data = environment_data)
##                 Df Variance        F Pr(>F)    
## Year_Collected   3 0.010548   9.3115  0.001 ***
## Shelf_Position   2 0.111391 147.5010  0.001 ***
## Residual       664 0.250722                    
## ---
## Signif. codes:  0 '***' 0.001 '**' 0.01 '*' 0.05 '.' 0.1 ' ' 1
\end{verbatim}

\begin{Shaded}
\begin{Highlighting}[]
\FunctionTok{sqrt}\NormalTok{(}\FunctionTok{vif.cca}\NormalTok{(my\_rda2))}
\end{Highlighting}
\end{Shaded}

\begin{verbatim}
##  Year_Collected2023  Year_Collected2024  Year_Collected2025   Shelf_PositionMid 
##            1.301675            1.303955            1.302811            1.143891 
## Shelf_PositionOuter 
##            1.134246
\end{verbatim}

\#\#\#Unexplained Variance Kaiser Guttman criterion

\begin{Shaded}
\begin{Highlighting}[]
\NormalTok{my\_rda2}\SpecialCharTok{$}\NormalTok{CA}\SpecialCharTok{$}\NormalTok{eig[my\_rda2}\SpecialCharTok{$}\NormalTok{CA}\SpecialCharTok{$}\NormalTok{eig }\SpecialCharTok{\textgreater{}} \FunctionTok{mean}\NormalTok{(my\_rda2}\SpecialCharTok{$}\NormalTok{CA}\SpecialCharTok{$}\NormalTok{eig)]}
\end{Highlighting}
\end{Shaded}

\begin{verbatim}
##        PC1        PC2 
## 0.09369417 0.07858813
\end{verbatim}

\section{Partial RDA}\label{partial-rda}

It is possible to run an RDA of a matrix Y, explained by a matrix
variables X, in the presence of co-variable(s) W.

This analysis allows to display the patterns of the response data (Y)
uniquely explained by a linear model of the explanatory variables (X)
when the effect of other covariates (W) is held constant (Borcard,
Gillet, and Legendre 2011). \#\# Hold Year constant Model the effect of
shelf position on coral recruit family abundances once the effect of
year is removed

\begin{Shaded}
\begin{Highlighting}[]
\NormalTok{partial\_rda }\OtherTok{\textless{}{-}} \FunctionTok{rda}\NormalTok{(recruit\_hellinger}\SpecialCharTok{\textasciitilde{}}\NormalTok{ Shelf\_Position }\SpecialCharTok{+} \FunctionTok{Condition}\NormalTok{(Year\_Collected), }\AttributeTok{data =}\NormalTok{ environment\_data)}
\end{Highlighting}
\end{Shaded}

\begin{Shaded}
\begin{Highlighting}[]
\FunctionTok{summary}\NormalTok{(partial\_rda)}
\end{Highlighting}
\end{Shaded}

\begin{verbatim}
## 
## Call:
## rda(formula = recruit_hellinger ~ Shelf_Position + Condition(Year_Collected),      data = environment_data) 
## 
## Partitioning of variance:
##               Inertia Proportion
## Total         0.37266     1.0000
## Conditioned   0.01055     0.0283
## Constrained   0.11139     0.2989
## Unconstrained 0.25072     0.6728
## 
## Eigenvalues, and their contribution to the variance 
## after removing the contribution of conditiniong variables
## 
## Importance of components:
##                         RDA1     RDA2     PC1     PC2     PC3    PC4
## Eigenvalue            0.1063 0.005075 0.09369 0.07859 0.04994 0.0285
## Proportion Explained  0.2936 0.014016 0.25874 0.21703 0.13792 0.0787
## Cumulative Proportion 0.2936 0.307614 0.56636 0.78338 0.92130 1.0000
## 
## Accumulated constrained eigenvalues
## Importance of components:
##                         RDA1     RDA2
## Eigenvalue            0.1063 0.005075
## Proportion Explained  0.9544 0.045563
## Cumulative Proportion 0.9544 1.000000
\end{verbatim}

Conditioned gives the amount of variance that has been explained by the
covariables and removed Constrained gives the amount of variance
uniquely explained by the explanatory variables

\begin{Shaded}
\begin{Highlighting}[]
\FunctionTok{RsquareAdj}\NormalTok{(partial\_rda)}
\end{Highlighting}
\end{Shaded}

\begin{verbatim}
## $r.squared
## [1] 0.2989068
## 
## $adj.r.squared
## [1] 0.2982176
\end{verbatim}

\subsection{Hold Shelf Position
Constant}\label{hold-shelf-position-constant}

Model the effect of year on coral recruit family abundances once the
effect of shelf position is removed

\begin{Shaded}
\begin{Highlighting}[]
\NormalTok{partial\_rda2 }\OtherTok{\textless{}{-}} \FunctionTok{rda}\NormalTok{(recruit\_hellinger}\SpecialCharTok{\textasciitilde{}}\NormalTok{ Year\_Collected }\SpecialCharTok{+} \FunctionTok{Condition}\NormalTok{(Shelf\_Position), }\AttributeTok{data =}\NormalTok{ environment\_data)}
\end{Highlighting}
\end{Shaded}

\begin{Shaded}
\begin{Highlighting}[]
\FunctionTok{summary}\NormalTok{(partial\_rda2)}
\end{Highlighting}
\end{Shaded}

\begin{verbatim}
## 
## Call:
## rda(formula = recruit_hellinger ~ Year_Collected + Condition(Shelf_Position),      data = environment_data) 
## 
## Partitioning of variance:
##                Inertia Proportion
## Total         0.372660    1.00000
## Conditioned   0.111980    0.30049
## Constrained   0.009958    0.02672
## Unconstrained 0.250722    0.67279
## 
## Eigenvalues, and their contribution to the variance 
## after removing the contribution of conditiniong variables
## 
## Importance of components:
##                           RDA1    RDA2     RDA3     PC1     PC2     PC3    PC4
## Eigenvalue            0.005869 0.00317 0.000919 0.09369 0.07859 0.04994 0.0285
## Proportion Explained  0.022513 0.01216 0.003525 0.35942 0.30147 0.19158 0.1093
## Cumulative Proportion 0.022513 0.03468 0.038200 0.39762 0.69910 0.89068 1.0000
## 
## Accumulated constrained eigenvalues
## Importance of components:
##                           RDA1    RDA2     RDA3
## Eigenvalue            0.005869 0.00317 0.000919
## Proportion Explained  0.589334 0.31838 0.092282
## Cumulative Proportion 0.589334 0.90772 1.000000
\end{verbatim}

\begin{Shaded}
\begin{Highlighting}[]
\FunctionTok{RsquareAdj}\NormalTok{(partial\_rda2)}
\end{Highlighting}
\end{Shaded}

\begin{verbatim}
## $r.squared
## [1] 0.02672155
## 
## $adj.r.squared
## [1] 0.02375285
\end{verbatim}

Test significance of partial RDA

\begin{Shaded}
\begin{Highlighting}[]
\FunctionTok{anova.cca}\NormalTok{(partial\_rda2)}
\end{Highlighting}
\end{Shaded}

\begin{verbatim}
## Permutation test for rda under reduced model
## Permutation: free
## Number of permutations: 999
## 
## Model: rda(formula = recruit_hellinger ~ Year_Collected + Condition(Shelf_Position), data = environment_data)
##           Df Variance      F Pr(>F)    
## Model      3 0.009958 8.7908  0.001 ***
## Residual 664 0.250722                  
## ---
## Signif. codes:  0 '***' 0.001 '**' 0.01 '*' 0.05 '.' 0.1 ' ' 1
\end{verbatim}

Since p≤0.05 we reject H0 and can conclude that the result was not
produced randomly.

\begin{Shaded}
\begin{Highlighting}[]
\FunctionTok{plot}\NormalTok{(partial\_rda2, }\AttributeTok{scaling =} \DecValTok{1}\NormalTok{)}
\end{Highlighting}
\end{Shaded}

\pandocbounded{\includegraphics[keepaspectratio]{Data_Analysis_Thesis_files/figure-latex/unnamed-chunk-38-1.pdf}}

\subsection{Variance partitioning}\label{variance-partitioning}

\begin{Shaded}
\begin{Highlighting}[]
\NormalTok{env\_shelf }\OtherTok{\textless{}{-}}\NormalTok{ environment\_data[,}\StringTok{"Shelf\_Position"}\NormalTok{]}

\NormalTok{env\_year }\OtherTok{\textless{}{-}}\NormalTok{ environment\_data[,}\StringTok{"Year\_Collected"}\NormalTok{]}
\NormalTok{vp }\OtherTok{\textless{}{-}} \FunctionTok{varpart}\NormalTok{(recruit\_hellinger, env\_shelf, env\_year)}
\NormalTok{vp}
\end{Highlighting}
\end{Shaded}

\begin{verbatim}
## 
## Partition of variance in RDA 
## 
## Call: varpart(Y = recruit_hellinger, X = env_shelf, env_year)
## 
## Explanatory tables:
## X1:  env_shelf
## X2:  env_year 
## 
## No. of explanatory tables: 2 
## Total variation (SS): 249.31 
##             Variance: 0.37266 
## No. of observations: 670 
## 
## Partition table:
##                      Df R.squared Adj.R.squared Testable
## [a+c] = X1            2   0.30049       0.29839     TRUE
## [b+c] = X2            3   0.02830       0.02393     TRUE
## [a+b+c] = X1+X2       5   0.32721       0.32214     TRUE
## Individual fractions                                    
## [a] = X1|X2           2                 0.29822     TRUE
## [b] = X2|X1           3                 0.02375     TRUE
## [c]                   0                 0.00017    FALSE
## [d] = Residuals                         0.67786    FALSE
## ---
## Use function 'rda' to test significance of fractions of interest
\end{verbatim}

\begin{Shaded}
\begin{Highlighting}[]
\FunctionTok{plot}\NormalTok{(vp)}
\end{Highlighting}
\end{Shaded}

\pandocbounded{\includegraphics[keepaspectratio]{Data_Analysis_Thesis_files/figure-latex/unnamed-chunk-39-1.pdf}}

\#\#\#Need to do test of significance of fractions

\begin{Shaded}
\begin{Highlighting}[]
\FunctionTok{anova.cca}\NormalTok{(}\FunctionTok{rda}\NormalTok{(recruit\_hellinger, env\_shelf))}
\end{Highlighting}
\end{Shaded}

\begin{verbatim}
## Permutation test for rda under reduced model
## Permutation: free
## Number of permutations: 999
## 
## Model: rda(X = recruit_hellinger, Y = env_shelf)
##           Df Variance      F Pr(>F)    
## Model      2  0.11198 143.26  0.001 ***
## Residual 667  0.26068                  
## ---
## Signif. codes:  0 '***' 0.001 '**' 0.01 '*' 0.05 '.' 0.1 ' ' 1
\end{verbatim}

\begin{Shaded}
\begin{Highlighting}[]
\FunctionTok{anova.cca}\NormalTok{(}\FunctionTok{rda}\NormalTok{(recruit\_hellinger, env\_year))}
\end{Highlighting}
\end{Shaded}

\begin{verbatim}
## Permutation test for rda under reduced model
## Permutation: free
## Number of permutations: 999
## 
## Model: rda(X = recruit_hellinger, Y = env_year)
##           Df Variance      F Pr(>F)    
## Model      3  0.01055 6.4666  0.001 ***
## Residual 666  0.36211                  
## ---
## Signif. codes:  0 '***' 0.001 '**' 0.01 '*' 0.05 '.' 0.1 ' ' 1
\end{verbatim}

\begin{Shaded}
\begin{Highlighting}[]
\CommentTok{\#[a + b+c]}
\FunctionTok{anova.cca}\NormalTok{(}\FunctionTok{rda}\NormalTok{(recruit\_hellinger, environment\_data))}
\end{Highlighting}
\end{Shaded}

\begin{verbatim}
## Permutation test for rda under reduced model
## Permutation: free
## Number of permutations: 999
## 
## Model: rda(X = recruit_hellinger, Y = environment_data)
##           Df Variance      F Pr(>F)    
## Model     20  0.16862 26.816  0.001 ***
## Residual 649  0.20404                  
## ---
## Signif. codes:  0 '***' 0.001 '**' 0.01 '*' 0.05 '.' 0.1 ' ' 1
\end{verbatim}

\begin{Shaded}
\begin{Highlighting}[]
\CommentTok{\# [a]}
\FunctionTok{anova.cca}\NormalTok{(}\FunctionTok{rda}\NormalTok{(recruit\_hellinger, env\_shelf, env\_year))}
\end{Highlighting}
\end{Shaded}

\begin{verbatim}
## Permutation test for rda under reduced model
## Permutation: free
## Number of permutations: 999
## 
## Model: rda(X = recruit_hellinger, Y = env_shelf, Z = env_year)
##           Df Variance     F Pr(>F)    
## Model      2  0.11139 147.5  0.001 ***
## Residual 664  0.25072                 
## ---
## Signif. codes:  0 '***' 0.001 '**' 0.01 '*' 0.05 '.' 0.1 ' ' 1
\end{verbatim}

\begin{Shaded}
\begin{Highlighting}[]
\CommentTok{\# [c]}
\FunctionTok{anova.cca}\NormalTok{(}\FunctionTok{rda}\NormalTok{(recruit\_hellinger, env\_year, env\_shelf))}
\end{Highlighting}
\end{Shaded}

\begin{verbatim}
## Permutation test for rda under reduced model
## Permutation: free
## Number of permutations: 999
## 
## Model: rda(X = recruit_hellinger, Y = env_year, Z = env_shelf)
##           Df Variance      F Pr(>F)    
## Model      3 0.009958 8.7908  0.001 ***
## Residual 664 0.250722                  
## ---
## Signif. codes:  0 '***' 0.001 '**' 0.01 '*' 0.05 '.' 0.1 ' ' 1
\end{verbatim}

\#Benthic Data

Importing benthic data

\begin{Shaded}
\begin{Highlighting}[]
\NormalTok{PIT\_50m }\OtherTok{\textless{}{-}} \FunctionTok{read\_csv}\NormalTok{(}\StringTok{"PIT\_50m.csv"}\NormalTok{) }\SpecialCharTok{\%\textgreater{}\%}
  \FunctionTok{mutate}\NormalTok{(}\FunctionTok{across}\NormalTok{(}\DecValTok{8}\SpecialCharTok{:}\DecValTok{50}\NormalTok{, }\SpecialCharTok{\textasciitilde{}}\FunctionTok{replace\_na}\NormalTok{(.x, }\DecValTok{0}\NormalTok{)))}\CommentTok{\# Read in data}

\NormalTok{PIT\_50m }\OtherTok{\textless{}{-}}\NormalTok{ PIT\_50m }\SpecialCharTok{\%\textgreater{}\%} 
  \FunctionTok{filter}\NormalTok{(Year }\SpecialCharTok{\%in\%} \FunctionTok{c}\NormalTok{(}\StringTok{"2021{-}22"}\NormalTok{, }\StringTok{"2022{-}23"}\NormalTok{, }\StringTok{"2023{-}24"}\NormalTok{, }\StringTok{"2024{-}25"}\NormalTok{, }\StringTok{"2025{-}26"}\NormalTok{))}

\NormalTok{Belt\_10m }\OtherTok{\textless{}{-}} \FunctionTok{read\_csv}\NormalTok{(}\StringTok{"Belt\_10m.csv"}\NormalTok{) }\SpecialCharTok{\%\textgreater{}\%}
  \FunctionTok{mutate}\NormalTok{(}\FunctionTok{across}\NormalTok{(}\DecValTok{8}\SpecialCharTok{:}\DecValTok{25}\NormalTok{, }\SpecialCharTok{\textasciitilde{}}\FunctionTok{replace\_na}\NormalTok{(.x, }\DecValTok{0}\NormalTok{))) }\CommentTok{\# Read in data}
\NormalTok{Belt\_10m }\OtherTok{\textless{}{-}}\NormalTok{ Belt\_10m }\SpecialCharTok{\%\textgreater{}\%}
  \FunctionTok{filter}\NormalTok{(Year }\SpecialCharTok{\%in\%} \FunctionTok{c}\NormalTok{(}\StringTok{"2021{-}22"}\NormalTok{, }\StringTok{"2022{-}23"}\NormalTok{, }\StringTok{"2023{-}24"}\NormalTok{, }\StringTok{"2024{-}25"}\NormalTok{, }\StringTok{"2025{-}26"}\NormalTok{))}
\end{Highlighting}
\end{Shaded}

\subsection{is the total number of recruits related to the number
proportion of coral in
PIT?}\label{is-the-total-number-of-recruits-related-to-the-number-proportion-of-coral-in-pit}

I am going to have to make a decision if I just keep the January or
October, and then will need to change the year to faciliate an inner
join.

\begin{Shaded}
\begin{Highlighting}[]
\CommentTok{\#PIT\_and\_recruits \textless{}{-} inner\_join(Coral\_settlement\_nGBR, PIT\_50m)}
\end{Highlighting}
\end{Shaded}

\section{Model Diagnostics Based on Murray
Feedback}\label{model-diagnostics-based-on-murray-feedback}

\begin{Shaded}
\begin{Highlighting}[]
\NormalTok{Coral\_settlement\_nGBR\_CS }\OtherTok{\textless{}{-}} \FunctionTok{read\_csv}\NormalTok{(}\StringTok{"\textasciitilde{}/Larval recruitment resources/Data Analysis/Coral{-}Settlement{-}Northern{-}GBR/Coral settlement\_nGBR\_CS.csv"}\NormalTok{) }\SpecialCharTok{\%\textgreater{}\%}
  \FunctionTok{mutate}\NormalTok{(}\FunctionTok{across}\NormalTok{(}\DecValTok{15}\SpecialCharTok{:}\DecValTok{29}\NormalTok{, }\SpecialCharTok{\textasciitilde{}}\FunctionTok{replace\_na}\NormalTok{(.x, }\DecValTok{0}\NormalTok{))) }\CommentTok{\# Read in data}

\NormalTok{Coral\_settlement\_nGBR }\OtherTok{\textless{}{-}}\NormalTok{ Coral\_settlement\_nGBR\_CS }\SpecialCharTok{\%\textgreater{}\%} 
  \FunctionTok{filter}\NormalTok{(}\StringTok{\textasciigrave{}}\AttributeTok{Shelf posn}\StringTok{\textasciigrave{}} \SpecialCharTok{!=} \StringTok{"Coral Sea"}\NormalTok{, }\CommentTok{\#remove coral sea data}
\NormalTok{        Obs }\SpecialCharTok{!=} \FunctionTok{c}\NormalTok{(}\DecValTok{1499}\NormalTok{, }\DecValTok{1621}\NormalTok{, }\DecValTok{1622}\NormalTok{)) }\CommentTok{\#Remove loose tiles with heavy biofouling}

\NormalTok{Coral\_settlement\_nGBR }\OtherTok{\textless{}{-}}\NormalTok{ Coral\_settlement\_nGBR }\SpecialCharTok{\%\textgreater{}\%} 
  \FunctionTok{rename}\NormalTok{(}\StringTok{"Shelf\_Position"}\OtherTok{=} \StringTok{\textasciigrave{}}\AttributeTok{Shelf posn}\StringTok{\textasciigrave{}}\NormalTok{,}
         \StringTok{"Year\_Collected"}\OtherTok{=} \StringTok{\textasciigrave{}}\AttributeTok{Year (Jan)}\StringTok{\textasciigrave{}}\NormalTok{,}
         \StringTok{"Site\_\#"}\OtherTok{=} \StringTok{\textasciigrave{}}\AttributeTok{Site \#}\StringTok{\textasciigrave{}}\NormalTok{,}
         \StringTok{"Acroporidae"} \OtherTok{=}\NormalTok{Acropora,}
         \StringTok{"Pocilloporidae"} \OtherTok{=}\NormalTok{ Pocillopora,}
         \StringTok{"Poritid"} \OtherTok{=}\NormalTok{ Porites) }\SpecialCharTok{|\textgreater{}} 
  \FunctionTok{mutate}\NormalTok{(}\AttributeTok{Year\_Collected =} \FunctionTok{factor}\NormalTok{(Year\_Collected),}
         \AttributeTok{Shelf\_Position =} \FunctionTok{factor}\NormalTok{(Shelf\_Position))}
\end{Highlighting}
\end{Shaded}

Suggests that the variance is not equal.

The 0 count skewed So we would do poisson or negative binomial

Priors would be beta, intercept, sd, and ``shape''?

\end{document}
